\section{Pairwise-Independent Hash (SODA Baseline)}
\label{sec:soda}

\subsection{The Hash Function}
\label{subsec:the-hash-function}

The paper's original locality-preserving hash~\citep[§3.1]{goswami2015approximate}:
\begin{equation}
    h(x) = \bigl(u(\lfloor x/2^K \rfloor) + (x \bmod 2^K)\bigr) \bmod 2^K
    \label{eq:soda-hash}
\end{equation}
That is: compute the block index $\lfloor x/2^K \rfloor$, hash it with $u$ to get a per-block rotation,
then add the key's within-block offset $x \bmod 2^K$.
The result is in $[0, 2^K)$.

$u: [0, U/2^K) \to [0, 2^K)$ is drawn from the multiply-add-shift family
of~\citet{dietzfelbinger1996universal}, stored as two 64-bit coefficients $(a, b)$:
\begin{equation}
    u(x) = \left\lfloor \frac{a \cdot x + b \bmod 2^{64}}{2^{64-K}} \right\rfloor.
    \label{eq:soda-u}
\end{equation}

\begin{figure}[h]
    \centering
    \includegraphics[width=0.85\textwidth]{soda_hash_mechanism}
    \caption{SODA hash mechanism: per-block rotation maps keys to $[0, 2^K)$.
    A query range spanning two blocks produces at most two intervals in
    the mapped space.}
    \label{fig:soda-mechanism}
\end{figure}

Properties of the hash:
\begin{itemize}
    \item \textbf{Pairwise-independent:} $\Pr_{h \sim \mathcal{H}}[h(x_1) = h(x_2)] \leq 2^{-K}$
    for every pair $x_1 \neq x_2 \in U$, where $\mathcal{H}$ is the multiply-add-shift family
    (Equation~\eqref{eq:soda-u}).
    This holds for sequential, clustered, or adversarial data.
    \item \textbf{Locality-preserving:} the hash applies a per-block rotation, so
    consecutive keys within the same block map to consecutive positions in $[0, 2^K)$.
    When $\mathcal{L} \leq 2^K$, a query range $[a, b]$ spans at most 2 blocks,
    producing at most 2 intervals in $[0, 2^K)$, each checked with one
    \textsc{IsEmpty} call to ERE.
    (The choice $K = \lceil\log_2(n\mathcal{L}/\varepsilon)\rceil \geq \log_2 \mathcal{L}$
    guarantees this.)
\end{itemize}

\paragraph{Comparison with Grafite.}
Grafite~\citep{grafite2024} uses a hash of the same locality-preserving form but draws
its per-block shift from a different pairwise-independent family --- multiply-add-modulo
over a Mersenne prime, requiring true modulo operations, where the multiply-add-shift
family of~\citet{dietzfelbinger1996universal} used here replaces them with a multiplication
and a right shift. The theoretical guarantees are identical;
Section~\ref{sec:grafite-deep-dive} dissects Grafite's full construction and our
implementation-level contributions to it.

\subsection{ARE Construction With SODA Hash}

This is a direct implementation of the construction from~\citet[§3]{goswami2015approximate}.
The hash design, its proofs of locality-preservation and pairwise independence, and the
resulting space/FPR bounds are all from the original paper --- our contribution is the
practical implementation.

Combining this hash with ERE (Chapter~\ref{chap:ere}) yields the baseline ARE filter:
\begin{enumerate}
    \item Divide $U$ into blocks of size $2^K$. Let $\text{blockIdx}(x) = \lfloor x / 2^K \rfloor$
    be the block index of key~$x$.
    \item For each block, compute pairwise-independent shift: top $K$ bits of
    $(a \cdot \text{blockIdx}(x) + b)$.
    \item Hash each key via Equation~\eqref{eq:soda-hash}.
    \item Sort, deduplicate, build ERE over $[0, 2^K)$.
    \item Query: hash both endpoints; rotation may split the range into two
    intervals --- check both in ERE.
\end{enumerate}

\paragraph{Filter properties.}
\begin{itemize}
    \item \textbf{FPR $\leq \varepsilon$ for any data distribution} --- follows from pairwise
    independence of the hash, with $K = \lceil \log_2(n\mathcal{L}/\varepsilon) \rceil$ giving
    BPK~$= \log_2(\mathcal{L}/\varepsilon) + O(1)$, matching the information-theoretic lower bound.
    Including ERE metadata, the total cost is $\log_2(\mathcal{L}/\varepsilon) + {\sim}2$ bits per key.
    \item \textbf{No false negatives} --- if key $x \in [a, b]$, then
    $h(x)$ lies in the mapped range $h([a, b])$, which is fully checked in ERE.
\end{itemize}

\subsection{Empirical Validation}
\label{subsec:empirical-validation}

On uniform data ($n = 2^{20}$, $\mathcal{L} = 128$), the measured FPR-vs-BPK curve
tracks the theoretical bound with a gap of ${\sim}2$~BPK (ERE overhead,
Figure~\ref{fig:soda-uniform}).

\begin{figure}[h]
    \centering
    \includegraphics[width=0.85\textwidth]{soda_tradeoff_uniform}
    \caption{FPR vs BPK on uniform keys ($n = 2^{20}$, $\mathcal{L} = 128$). SODA tracks
    the theoretical bound with ${\sim}2$~BPK overhead from ERE metadata.}
    \label{fig:soda-uniform}
\end{figure}
